% ----- CHAUPAI 34 -----

\newpage

\subsection*{\deva{चतुस्त्रिंशत्तम चौपाई} (Thirty-Fourth Chaupai)}
\addcontentsline{toc}{subsection}{\deva{चतुस्त्रिंशत्तम चौपाई} (Thirty-Fourth Chaupai) :   \deva{अन्त काल...}}

\begin{minipage}[t]{0.55\textwidth}
\vspace{0pt}

{\large \linespread{1.5}\selectfont
\deva{अंत काल रघुबर पुर जाई।}\\
\deva{जहाँ जन्म हरि-भक्त कहाई॥}
\par}

\vspace{0.5em}

{\large \linespread{1.5}\selectfont
\telu{అంత కాల రఘుబర పుర జాయీ।}\\
\telu{జహాఁ జన్మ హరి-భక్త కహాయీ॥}
\par}

\vspace{0.5em}

{\large \linespread{1.3}\selectfont
\textit{anta kāla raghubara pura jāī |}\\
\textit{jahā̃ janma hari-bhakta kahāī ||}
\par}
\end{minipage}%
\hfill
\begin{minipage}[t]{0.42\textwidth}
\vspace{0pt}
\begin{tcolorbox}[anchorbox]
\textbf{The Ultimate Destination:} By dedicating our lives to service and excellence, our legacy (\textit{Anta Kaal}) becomes immortal. In the *Valmiki Ramayana* (*Yuddha Kanda, 1:13*), Lord Rama laments that Hanuman's monumental, selfless service is so vast that no material reward could ever match it. Rama offers his own embrace as his entire wealth, establishing Hanuman's eternal spiritual legacy.
\vspace{0.5em}\newline Rama's ultimate embrace and praise in Yuddha Kanda established Hanuman's immortal, spiritual legacy above all material rewards.\newline\null\hfill\href{https://www.valmiki.iitk.ac.in/sloka?field_kanda_tid=6&language=dv&field_sarga_value=1}{\textit{Yuddha Kanda, 1:13}}
\end{tcolorbox}
\end{minipage}

\vspace{0.5em}
\subsubsection*{\textbf{Visual Rhythm Map:}}
\begin{table}[h!]
\centering
\renewcommand{\arraystretch}{1.2}
\setlength{\tabcolsep}{0pt}
\begin{tabularx}{\textwidth}{|*{16}{>{\centering\arraybackslash\hsize=1.0\hsize}X|}}
\hline
\multicolumn{3}{|>{\centering\arraybackslash\hsize=3.0\hsize}X|}{\textbf{\guru{अं}\laghu{त}} \par (3)} &
\multicolumn{3}{>{\centering\arraybackslash\hsize=3.0\hsize}X|}{\textbf{\guru{का}\laghu{ल}} \par (3)} &
\multicolumn{4}{>{\centering\arraybackslash\hsize=4.0\hsize}X|}{\textbf{\laghu{र}\laghu{घु}\laghu{ब}\laghu{र}} \par (4)} &
\multicolumn{2}{>{\centering\arraybackslash\hsize=2.0\hsize}X|}{\textbf{\laghu{पु}\laghu{र}} \par (2)} &
\multicolumn{4}{>{\centering\arraybackslash\hsize=4.0\hsize}X|}{\textbf{\guru{जा}\guru{ई}} \par (4)} \\
\hline
\end{tabularx}
\vspace{-1.5em}
\end{table}

\begin{table}[h!]
\centering
\renewcommand{\arraystretch}{1.2}
\setlength{\tabcolsep}{0pt}
\begin{tabularx}{\textwidth}{|*{16}{>{\centering\arraybackslash\hsize=1.0\hsize}X|}}
\hline
\multicolumn{3}{|>{\centering\arraybackslash\hsize=3.0\hsize}X|}{\textbf{\laghu{ज}\guru{हाँ}} \par (3)} &
\multicolumn{3}{>{\centering\arraybackslash\hsize=3.0\hsize}X|}{\textbf{\guru{ज}\laghu{न्म}} \par (3)} &
\multicolumn{5}{>{\centering\arraybackslash\hsize=5.0\hsize}X|}{\textbf{\laghu{ह}\laghu{रि}\guru{भ}\laghu{क्त}} \par (5)} &
\multicolumn{5}{>{\centering\arraybackslash\hsize=5.0\hsize}X|}{\textbf{\laghu{क}\guru{हा}\guru{ई}} \par (5)} \\
\hline
\end{tabularx}
\vspace{-1.5em}
\end{table}

\vspace{1em}
\textbf{ \deva{सरल-संस्कृतम्}:}\\
 \deva{सः अन्त्यकाले श्रीरामस्य धाम (वैकुण्ठं) गच्छति।}\\
 \deva{यत्र तस्य जन्म भवति चेत्, सः भगवतः भक्तः इति उच्यते॥}\\


At the end of time/life, one goes to the celestial abode of Lord Rama.\\
And wherever they might be born anew, they are renowned as a devotee of Hari/God.


\vspace{1em}
\newpage
\textbf{\deva{प्रतिपदार्थः} (Meaning of each word):}
\begin{table}[h!]
\centering
\renewcommand{\arraystretch}{1.2}
\begin{tabularx}{\textwidth}{@{} l l X X @{}}
\toprule
\textbf{Awadhi} & \textbf{Sanskrit} & \textbf{English} & \textbf{Grammar \& Notes} \\
\midrule
\deva{अन्त काल} / \telu{అన్త కాల} / \textit{anta kāla}  &  \deva{अन्त्य-काले}  &  At the end of life  &   \\
\deva{रघुबर पुर} / \telu{రఘుబర పుర} / \textit{raghubara pura}  &  \deva{श्रीराम-धाम}  &  Rama's celestial city  & \\ 
\deva{जाई} / \telu{జాయీ} / \textit{jāī}  &  \deva{गच्छति}  &  Goes / Reaches  &   \\
\deva{जहाँ} / \telu{జహాఁ} / \textit{jahā̃}  &  \deva{यत्र}  &  Where / Wherever  &   \\
\deva{जन्म} / \telu{జన్మ} / \textit{janma}  &  \deva{जन्म}  &  Birth  & Pure Sanskrit \\
\deva{हरिभक्त} / \telu{హరిభక్త} / \textit{hari-bhakta}  &  \deva{हरि-भक्तः}  &  Devotee of God  &   \\
\deva{कहाई} / \telu{కహాయీ} / \textit{kahāī}  &  \deva{कथ्यते / उच्यते}  &  Is called / renowned  & Passive voice \\
\bottomrule
\end{tabularx}
\end{table}

\begin{tcolorbox}[poetrybox]
\textbf{Meter Alert (Heavy Conjuncts):}\\
Pay attention to \deva{जन्म} (Janma) and \deva{हरिभक्त} (Haribhakta) in Line 2. Because of the `nm` and `kt` conjuncts, the preceding vowels (\deva{ज} and \deva{भ}) absorb an extra beat, making them heavy (Guru) to maintain the exact 16-beat rhythm.
\end{tcolorbox}

\newpage
